% Detailed solutions to the 18 items of Listas/pratica-lecture-03.md (Lecture 3: mixed strategies
% and dynamic games). Item order: sol-items.tex. Eleven items are reused from
% ../lista-reais-lectures-01-05-solutions/ (\solitemr); seven are new here (\solitem).
% Compile here: pdflatex main.tex (twice). Check: pdffonts main.pdf (all Type 1, uni = yes).
\documentclass[11pt,a4paper]{article}
\newcommand{\soltema}{Practice list, Lecture 3: detailed solutions}
% =============================================================================
% Shared preamble of the solutions to the real-exercise list for Lectures 1--5.
% Input by main.tex (all lectures) and lecture-0N.tex (one lecture). The driver defines
% \soltema (running head) before \input-ing this file. Layout copied from
% Listas/lista-reais-lectures-01-04-solutions/main.tex.
% Compile from this folder: pdflatex <driver>.tex (twice); pdffonts <driver>.pdf (all Type 1).
% =============================================================================
\usepackage[utf8]{inputenc}
\usepackage[T1]{fontenc}
\usepackage{lmodern}   % Type 1 vector fonts (replaces PK bitmaps)
\usepackage{cmap}      % ToUnicode CMap -> copyable text and accents
\usepackage[english]{babel}
\usepackage{amsmath,amssymb,amsthm,mathtools}
\usepackage[a4paper,margin=2.4cm]{geometry}
\usepackage{microtype}
\usepackage{booktabs,tabularx,array}
\usepackage{enumitem}
\usepackage{xcolor}
\usepackage{tikz}
\usetikzlibrary{arrows.meta,calc}
\usepackage[most]{tcolorbox}
\usepackage{fancyhdr}
\usepackage[hidelinks]{hyperref}

\setlist{topsep=3pt,itemsep=2pt}
\setlength{\parindent}{0pt}
\setlength{\parskip}{5pt}
\renewcommand{\arraystretch}{1.2}

\pagestyle{fancy}\fancyhf{}
\lhead{\footnotesize Economics of Organizations and Incentives --- MPE Insper}
\rhead{\footnotesize \soltema}
\cfoot{\footnotesize\thepage}
\renewcommand{\headrulewidth}{0.3pt}
\setlength{\headheight}{14pt}

\definecolor{hdr}{RGB}{30,80,150}
\definecolor{ans}{RGB}{20,110,70}

% ----------------------------------------------------------------------------- numbering
% Item files carry no numbers. A driver calls \solitem{<combined n>}{<slug>}: it steps the
% per-lecture counter ItemNo, sets \CombNo and inputs sol-<slug>.tex (skipped with a log
% message if the file does not exist yet). \sollecture{<title>} opens a lecture and resets
% ItemNo, so items are numbered 1--k within each lecture in every driver.
\newcounter{ItemNo}
\newcommand{\CombNo}{0}
\newcommand{\solitem}[2]{\stepcounter{ItemNo}\renewcommand{\CombNo}{#1}%
  \IfFileExists{sol-#2.tex}{\input{sol-#2}}{\typeout{solitem: sol-#2.tex missing, skipped}}}
% The page break sits here, so a lecture's first item follows its heading on the same page.
\newif\ifafterlecture
\newcommand{\sollecture}[1]{\clearpage\phantomsection\part*{#1}\addcontentsline{toc}{part}{#1}%
  \setcounter{ItemNo}{0}\afterlecturetrue}

% \begin{exsol}{<source and exercise>}{<printed page>}{<level>}{<difficulty 1-5>} ... \end{exsol}
% One exercise of the list; the heading shows its number in the lecture and in the combined list.
\newenvironment{exsol}[4]{%
  \ifafterlecture\global\afterlecturefalse\else\clearpage\fi\phantomsection
  \addcontentsline{toc}{section}{Item \theItemNo\ (\CombNo): #1 (p.~#2)}%
  \section*{\large Item \theItemNo\ {\normalsize\mdseries(item \CombNo\ of the combined list)}:
    #1\hfill{\normalsize\mdseries #3 $\cdot$ difficulty #4/5}}%
  \textbf{Source:} #1, page~#2. The statement is reproduced verbatim in the list (item
  \CombNo); it is not repeated here.\par}{\par}
% Fixed sub-heads, in this order inside each exsol.
\newcommand{\asked}{\paragraph{(a) What is asked.}}
\newcommand{\method}{\paragraph{(b) Method.}}
\newcommand{\solution}{\paragraph{(c) Solution.}}
\newcommand{\interp}{\paragraph{(d) Interpretation.}}
\newcommand{\verify}{\paragraph{(e) Verification.}}
\newcommand{\mistakes}{\paragraph{(f) Common mistakes.}}
% Boxed final answer: \answer{...}
\newcommand{\answer}[1]{\begin{tcolorbox}[colback=ans!5,colframe=ans,boxrule=0.6pt,
  left=5pt,right=5pt,top=3pt,bottom=3pt]\textbf{Answer.} #1\end{tcolorbox}}
% Best-response mark in a payoff matrix: \br{8} = payoff 8 is a best reply of its player.
\newcommand{\br}[1]{\underline{#1}}
\newcommand{\MG}{Mu\~noz-Garcia}
\newcommand{\Sm}{S_{-i}}
\newcommand{\sm}{s_{-i}}

% Item order of the Lecture 3 solutions, as in Listas/pratica-lecture-03.md.
% \solitem{n}{slug}  : new solution, file sol-<slug>.tex in this folder.
% \solitemr{n}{slug} : reused solution, file in ../lista-reais-lectures-01-05-solutions/.
\newcommand{\solitemr}[2]{\stepcounter{ItemNo}\renewcommand{\CombNo}{#1}%
  \input{../lista-reais-lectures-01-05-solutions/sol-#2}}
\newcommand{\lectureiiititle}{Lecture 3 --- Mixed strategies and dynamic games}
\newcommand{\itemsiii}{%
\solitem{1}{tadelis-6-5}
\solitemr{2}{tadelis-7-1}
\solitemr{3}{munoz-garcia-4-3}
\solitemr{4}{tadelis-7-2}
\solitem{5}{munoz-garcia-4-4}
\solitemr{6}{tadelis-7-6}
\solitem{7}{munoz-garcia-4-9}
\solitemr{8}{tadelis-7-4}
\solitemr{9}{tadelis-7-9}
\solitemr{10}{tadelis-7-10}
\solitemr{11}{tadelis-8-5}
\solitemr{12}{tadelis-8-10}
\solitemr{13}{tadelis-8-11}
\solitem{14}{gibbons-1-14}
\solitem{15}{tadelis-6-3}
\solitemr{16}{mwg-9-b-4}
\solitemr{17}{gibbons-2-1}
\solitemr{18}{mwg-9-b-5}}


% The shared exsol heading says "of the combined list"; here the item number is the list's own.
\renewenvironment{exsol}[4]{%
  \ifafterlecture\global\afterlecturefalse\else\clearpage\fi\phantomsection
  \addcontentsline{toc}{section}{Item \theItemNo: #1 (p.~#2)}%
  \section*{\large Item \theItemNo: #1\hfill{\normalsize\mdseries #3 $\cdot$ difficulty #4/5}}%
  \textbf{Source:} #1, page~#2. The statement is reproduced verbatim in
  \texttt{Listas/pratica-lecture-03.md} (item \theItemNo); it is not repeated here.\par}{\par}

\title{\textbf{Practice List, Lecture 3: Detailed Solutions}\\[4pt]
  \large Economics of Organizations and Incentives (MPE Insper)}
\author{Solutions to the 18 exercises of \texttt{Listas/pratica-lecture-03.md}}
\date{November 2026}

% Prints each exercise's statement, verbatim from its practice list, right before its solution.
% Input from a solutions main.tex just before \begin{document}. The fragments stmt/item-NN.tex are
% generated by Listas/make_statements.py from Listas/pratica-lecture-0N.md (the single source).
% Assumes the counter ItemNo (sol-preamble.tex) and the switch \ifafterlecture exist.
\usepackage{longtable,booktabs,array,calc,graphicx}
\makeatletter
% Support macros that pandoc's LaTeX output expects.
\providecommand{\tightlist}{\setlength{\itemsep}{0pt}\setlength{\parskip}{0pt}}
\@ifundefined{c@none}{\newcounter{none}}{}  % pandoc's uncaptioned tables use \LTcaptype{none}
\providecommand{\pandocbounded}[1]{\begingroup
  \setkeys{Gin}{width=0.75\linewidth,height=0.32\textheight,keepaspectratio}#1\endgroup}
\definecolor{stmtrule}{RGB}{30,80,150}
% \statementboxn{<item number>}: the statement between two rules; warns if the fragment is missing.
\newcommand{\statementboxn}[1]{%
  \edef\stmt@file{stmt/item-\two@digits{#1}.tex}%
  \IfFileExists{\stmt@file}{%
    \par\smallskip\noindent{\color{stmtrule}\rule{\linewidth}{0.6pt}}\par\nopagebreak
    \noindent{\color{stmtrule}\textbf{Statement}} {\small(verbatim from the list, item~#1)}\par
    \nopagebreak{\small\input{\stmt@file}}\par
    \noindent{\color{stmtrule}\rule{\linewidth}{0.6pt}}\par\medskip}%
    {\PackageWarning{statement-hook}{No statement file for item #1}}}
\makeatother
% The 4-argument exsol of sol-preamble.tex (Lectures 2--8), with the statement after the heading.
\renewenvironment{exsol}[4]{%
  \ifafterlecture\global\afterlecturefalse\else\clearpage\fi\phantomsection
  \addcontentsline{toc}{section}{Item \theItemNo: #1 (p.~#2)}%
  \section*{\large Item \theItemNo: #1\hfill{\normalsize\mdseries #3 $\cdot$ difficulty #4/5}}%
  \statementboxn{\arabic{ItemNo}}}{\par}
  % statement of each item before its solution

\begin{document}
\maketitle
\thispagestyle{empty}

\noindent\sloppy These are detailed solutions to the 18 book exercises of the list
\texttt{pratica-lecture-03.md} (Lecture~3: mixed strategies, extensive form, backward induction and
subgame perfection), in the list's order from easy to hard; each heading shows the item's number,
level and difficulty from 1 (apply a definition) to 5 (graduate level). Each statement is printed verbatim from the list
right before its solution. Every solution has the same six parts: (a) what is asked; (b) the method, with
book section, page and lecture slide; (c) the solution, every step shown; (d) the economic
interpretation; (e) the verification, by \texttt{check\_L03.py} in this folder (new items) or the check
scripts of \texttt{lista-reais-lectures-01-05-solutions/} (reused items), and against the book's worked
answer where one exists; (f) common mistakes. Slide numbers in the reused solutions refer to the
drafts of the earlier real-exercise list and may differ by a few slides from the final
\texttt{Aula/Lecture03.pdf}. In bimatrices player~1 picks rows and the first payoff is player~1's;
\underline{underlined} payoffs are best replies.

\tableofcontents

\sollecture{\lectureiiititle}
\itemsiii

\end{document}
